\begin{lemma}[Support respects presentation equivalence]
If two raw presentations are structurally equivalent, then they have the same
support:
\[
\mathsf{PowerPresentationEq}(X,Y)
\quad\Longrightarrow\quad
\mathsf{PowerSupport}(X)=\mathsf{PowerSupport}(Y).
\]
\end{lemma}
\begin{proof}
Use the recursive definition of structural equivalence from
\noderef{PowerPresentationEq}.  In the atom case the equivalence says that
the two atom values are equal, and the atom support equation in
\noderef{PowerSupportEquations} gives equal singleton supports.  The mixed
atom--node cases are impossible by the disjointness characterization in
\noderef{PowerPresentationEqShape}.  In the node case, prove mutual support
inclusion.  If an atom belongs to the support of the first node, the node
support equation in \noderef{PowerSupportEquations} places it in the support
of some child $f(i)$.  The forward matching clause supplies $j$ with
$\mathsf{PowerPresentationEq}(f(i),g(j))$; the induction hypothesis places
the atom in the support of $g(j)$, hence in the support of the second node.
The backward matching clause gives the reverse inclusion.  Extensionality of
sets now yields equality of supports.
\end{proof}
